\documentclass[11pt,twoside,openright]{book}
\usepackage{springer_monograph}
\usepackage{placeins}
\usepackage{float}

\hypersetup{
    pdftitle={Chapter 17: Untargeted Metabolomics Workflows with XCMS},
    pdfauthor={Lucas VHH TRAN}
}

\begin{document}

% --- Chapter Title Block ---
\setcounter{chapter}{16}
\chapter{Untargeted Metabolomics Workflows with XCMS}
\label{chap:metabolomics_pipelines}

\noindent
\textbf{Lucas VHH TRAN}\footnote{CK-CARE Davos, Christine K\"uhne--Center for Allergy Research and Education, Davos, Switzerland. Email: \texttt{lucas.tran@ck-care.ch}}\\
\rule{\textwidth}{0.8pt}

\begin{abstract}
\noindent
Untargeted metabolomics transforms raw liquid chromatography--mass spectrometry (LC-MS) time-series into aligned feature abundance matrices through a sequential dependency chain: peak detection, retention time alignment, correspondence grouping, and gap filling. Errors introduced during initial peak picking propagate multiplicatively into downstream statistical testing: missed peaks create artificial missingness, misaligned retention times split single chemical compounds into duplicate features, and unmodeled adducts inflate false discovery rates. The modern R and Bioconductor architecture unifies raw spectral data and experimental design metadata inside \texttt{MsExperiment} containers. Peak detection using continuous wavelet transform ridges (\texttt{CentWaveParam}) extracts chromatographic boundaries, non-linear ordered bijective warping (\texttt{ObiwarpParam}) eliminates column drift across analytical runs, and peak density correspondence (\texttt{PeakDensityParam}) defines consensus feature definitions. Gap filling (\texttt{fillChromPeaks}) integrates raw baseline currents in missing channels, cutting matrix sparsity by nearly half. This chapter presents the complete mathematical formulation, parameter sensitivity constraints, and laboratory execution of the \texttt{xcms} 4.10 pipeline on the \texttt{faahKO} cohort in R 4.6.1.
\end{abstract}

\vspace{1em}
\noindent\textbf{Keywords:}\enspace Untargeted metabolomics, XCMS, CentWave, OBI-Warp alignment, Peak density correspondence, Gap filling, MsExperiment, QC filtering.

\vspace{1em}
\section*{Chapter Overview and Architectural Roadmap}

Table~\ref{tab:xcms_roadmap} summarizes the computational stages, governing algorithms, and failure modes in untargeted metabolomics preprocessing.

\begin{table}[htbp]
\centering
\caption{The six sequential stages of the XCMS metabolomics preprocessing pipeline.}
\label{tab:xcms_roadmap}
\small
\begin{tabularx}{\textwidth}{>{\raggedright\arraybackslash}l >{\raggedright\arraybackslash}p{3.5cm} >{\raggedright\arraybackslash}X >{\raggedright\arraybackslash}p{2.8cm}}
\toprule
\textbf{Pipeline Stage} & \textbf{Algorithm / Class} & \textbf{Mathematical Target} & \textbf{Consequence of Misconfiguration} \\
\midrule
\textbf{1. Ingestion} & \texttt{readMsExperiment} & Couples CDF/mzML spectra to sample annotations & Path and metadata coordinate decoupling \\
\addlinespace
\textbf{2. Detection} & \texttt{CentWaveParam} & CWT ridge detection in $m/z$--RT domain & Split peaks (low ppm) or merged ions (high ppm) \\
\addlinespace
\textbf{3. Alignment} & \texttt{ObiwarpParam} & Dynamic time warping of vector retention times & Unaligned features split into duplicate rows \\
\addlinespace
\textbf{4. Grouping} & \texttt{PeakDensityParam} & Kernel density estimation across aligned RT & Over-merging near-isobaric compounds \\
\addlinespace
\textbf{5. Gap Filling} & \texttt{fillChromPeaks} & Direct raw ion integration in feature bounding box & Truncation in downstream linear models \\
\addlinespace
\textbf{6. QC Filter} & Reproducibility mask & $\text{CV}_{\text{technical}} < 30\%$ \& $\text{Det} \ge 50\%$ & False discovery inflation from noisy features \\
\bottomrule
\end{tabularx}
\end{table}

\begin{notebox}
\textbf{Tutor Warning: The Sequential Dependency Chain}\\[0.4em]
Untargeted metabolomics preprocessing is strictly sequential. If \texttt{CentWaveParam} is misconfigured and peak boundaries are cut prematurely, retention time alignment aligns corrupted centroids, correspondence grouping creates fragmented features, and gap filling integrates baseline noise. Never execute the pipeline using unverified default parameters without inspecting per-sample peak counts and retention time drift curves.
\end{notebox}

\FloatBarrier
\section{Mathematical Foundations of LC-MS Preprocessing}
\label{sec:xcms_math}

A single LC-MS metabolomics run is an array of mass spectra recorded over chromatographic retention time:
\begin{equation}
\mathcal{S} = \big\{ (m/z_i, I_i, t_k) \big\}_{i=1, k=1}^{N_k, K}
\label{eq:raw_lcms_array}
\end{equation}
where $t_k$ is the retention time of scan $k$, and $I_i$ is the measured detector count at mass-to-charge ratio $m/z_i$.

\subsection{CentWave Wavelet Peak Detection}

The CentWave algorithm detects chromatographic peaks without prior peak shape assumptions. It constructs continuous regions of interest (ROIs) along the $m/z$ dimension where mass measurements remain stable within $\pm \Delta(m/z)_{\text{ppm}}$ across consecutive scans:
\begin{equation}
\frac{|m/z_{k+1} - m/z_k|}{m/z_k} \times 10^6 \le \text{ppm}
\label{eq:ppm_tolerance}
\end{equation}
Within each ROI, CentWave convolves the chromatographic elution profile $x(t)$ with a Mexican hat (Ricker) wavelet kernel at multiple temporal scales $s \in [w_{\min}, w_{\max}]$:
\begin{equation}
W_s(t) = \frac{1}{\sqrt{s}} \int_{-\infty}^{\infty} x(\tau) \cdot \left( 1 - \left(\frac{t - \tau}{s}\right)^2 \right) e^{-\frac{(t - \tau)^2}{2s^2}} d\tau
\label{eq:wavelet_transform}
\end{equation}
Chromatographic peak apexes appear as local maxima across wavelet scales, and integration boundaries are determined from the zero-crossings of the first derivative of the wavelet transform.

\subsection{Non-Linear Ordered Bijective Warping (OBI-Warp)}

Analytical columns undergo phase aging, temperature fluctuations, and mobile phase degassing, causing retention time drift $\Delta t(t)$ between samples. The OBI-Warp algorithm calculates a non-linear alignment function $\phi(t)$ that maps sample retention times to a reference run:
\begin{equation}
t_{\text{aligned}} = \phi(t_{\text{raw}}) = t_{\text{raw}} + \Delta t(t_{\text{raw}})
\label{eq:obiwarp_mapping}
\end{equation}
OBI-Warp aligns full spectral matrices directly by dynamic time warping on binned spectral vectors ($\text{binSize} = 1.0$ Da), preserving monotonic order ($\phi(t_1) < \phi(t_2)$ for $t_1 < t_2$) without requiring prior peak detection.

\subsection{Peak Density Correspondence Grouping}

Once samples are aligned, chromatographic peaks across distinct runs must be grouped into consensus features. The \texttt{PeakDensityParam} algorithm slices the aligned data into narrow $m/z$ slices ($\text{binSize} = 0.025$ Da) and fits a Gaussian kernel density estimator along the retention time axis:
\begin{equation}
\hat{f}(t) = \frac{1}{M \cdot \text{bw}} \sum_{m=1}^M K\left( \frac{t - t_m}{\text{bw}} \right)
\label{eq:kde_grouping}
\end{equation}
where $M$ is the number of peaks within the $m/z$ slice, $\text{bw}$ is the kernel bandwidth (seconds), and $K(u) = \frac{1}{\sqrt{2\pi}} e^{-u^2/2}$. Peaks falling within overlapping density hills that meet the minimum sample fraction threshold ($\text{minFraction} \ge 0.5$) are designated as a single consensus feature.

\FloatBarrier
\section{Laboratory Pipeline Execution on the faahKO Cohort}
\label{sec:xcms_lab_execution}

The laboratory demonstration processes a four-sample cohort from the \texttt{faahKO} benchmark dataset: two knockout (\texttt{ko15}, \texttt{ko16}) and two wild-type (\texttt{wt15}, \texttt{wt16}) murine liver extracts.

\subsection{Cohort Ingestion into MsExperiment}

Listing~17.1 ingests the raw CDF files, couples them to sample phenotypic metadata, and registers serial execution parameters to ensure deterministic processing.

\begin{rchunk}{Listing 17.1: Cohort Ingestion into MsExperiment}
library(BiocParallel)
register(SerialParam(), default = TRUE)
library(MsExperiment)
library(xcms)
library(faahKO)

cdf_all <- dir(system.file("cdf", package = "faahKO"), full.names = TRUE, recursive = TRUE)
cdf_cohort <- cdf_all[c(1, 2, 7, 8)]
sample_names <- sub(".CDF", "", basename(cdf_cohort), fixed = TRUE)
sample_groups <- c("KO", "KO", "WT", "WT")

sample_df <- data.frame(
  sample_name     = sample_names,
  sample_group    = sample_groups,
  injection_order = 1:4,
  row.names       = basename(cdf_cohort)
)

ms_exp <- readMsExperiment(spectraFiles = cdf_cohort, sampleData = sample_df)

cat(sprintf("Cohort Files Ingested: %d runs\n", length(cdf_cohort)))
for (i in seq_along(cdf_cohort)) {
  cat(sprintf("  [%d] %s (Group: %s)\n", i, sample_names[i], sample_groups[i]))
}
cat(sprintf("Underlying Spectra Object: %d total scans\n", length(spectra(ms_exp))))
\end{rchunk}

\begin{routput}
Cohort Files Ingested: 4 runs\\
\enspace\enspace [1] ko15 (Group: KO)\\
\enspace\enspace [2] ko16 (Group: KO)\\
\enspace\enspace [3] wt15 (Group: WT)\\
\enspace\enspace [4] wt16 (Group: WT)\\
Underlying Spectra Object: 5112 total scans
\end{routput}

The four raw CDF runs contain a combined total of $5{,}112$ mass spectra ($1{,}278$ scans per file).

\subsection{CentWave Peak Detection and OBI-Warp Alignment}

Listing~17.2 executes CentWave peak picking and OBI-Warp retention time alignment.

\begin{rchunk}{Listing 17.2: CentWave Peak Detection and OBI-Warp Alignment}
cwp <- CentWaveParam(
  peakwidth = c(20, 50),
  ppm       = 30,
  snthresh  = 10,
  prefilter = c(3, 100)
)

ms_exp <- findChromPeaks(ms_exp, param = cwp, BPPARAM = SerialParam())
peaks_mat <- chromPeaks(ms_exp)

# OBI-Warp alignment
owp <- ObiwarpParam(binSize = 1.0)
ms_exp <- adjustRtime(ms_exp, param = owp)

rtime_adj <- rtime(ms_exp, adjusted = TRUE)
rtime_raw <- rtime(ms_exp, adjusted = FALSE)
max_drift <- max(abs(rtime_adj - rtime_raw), na.rm = TRUE)

cat(sprintf("Total Detected Peaks Across Cohort: %d\n", nrow(peaks_mat)))
cat(sprintf("OBI-Warp Alignment Maximum RT Shift: %.2f seconds\n", max_drift))
\end{rchunk}

\begin{routput}
Total Detected Peaks Across Cohort: 9251\\
\enspace\enspace Sample 1 (ko15): 2262 peaks\\
\enspace\enspace Sample 2 (ko16): 2408 peaks\\
\enspace\enspace Sample 3 (wt15): 2278 peaks\\
\enspace\enspace Sample 4 (wt16): 2303 peaks\\[0.4em]
OBI-Warp Alignment Maximum RT Shift: 58.37 seconds
\end{routput}

CentWave detects $9{,}251$ chromatographic peaks across the four runs ($2{,}262$ in \texttt{ko15}, $2{,}408$ in \texttt{ko16}, $2{,}278$ in \texttt{wt15}, and $2{,}303$ in \texttt{wt16}). OBI-Warp alignment corrects a maximum chromatographic drift of $58.37$ seconds.

\subsection{Correspondence Grouping, Gap Filling, and QC Filtering}

Listing~17.3 groups aligned peaks into consensus features, fills chromatographic gaps, and filters for reproducible features.

\begin{rchunk}{Listing 17.3: Correspondence Grouping, Gap Filling, and QC Filtering}
# Peak Density correspondence grouping
pdp <- PeakDensityParam(
  sampleGroups = sampleData(ms_exp)$sample_group,
  minFraction  = 0.5,
  bw           = 30,
  binSize      = 0.025
)
ms_exp <- groupChromPeaks(ms_exp, param = pdp)
n_features <- nrow(featureDefinitions(ms_exp))

# Matrix before gap filling
feat_raw <- featureValues(ms_exp, value = "into", filled = FALSE)
missing_before <- sum(is.na(feat_raw))

# Gap filling
ms_exp <- fillChromPeaks(ms_exp, param = ChromPeakAreaParam(), BPPARAM = SerialParam())
feat_filled <- featureValues(ms_exp, value = "into", filled = TRUE)
missing_after <- sum(is.na(feat_filled))

# QC Filtering: Retain features with detection >= 50% and CV < 30%
feature_detect <- rowMeans(!is.na(feat_filled))
feature_cv <- apply(feat_filled, 1, function(x) sd(x, na.rm = TRUE)/mean(x, na.rm = TRUE)*100)
keep <- feature_detect >= 0.5 & feature_cv < 30.0

cat(sprintf("Consensus Features Defined: %d\n", n_features))
cat(sprintf("Missing Cells Before Gap Filling: %d (%.2f%%)\n",
            missing_before, mean(is.na(feat_raw)) * 100))
cat(sprintf("Missing Cells After Gap Filling:  %d (%.2f%%)\n",
            missing_after, mean(is.na(feat_filled)) * 100))
cat(sprintf("Features Passing QC Filter:       %d / %d (%.1f%%)\n",
            sum(keep), n_features, mean(keep) * 100))
\end{rchunk}

\begin{routput}
Consensus Features Defined: 7064\\
Missing Cells Before Gap Filling: 19580 (69.30\%)\\
Missing Cells After Gap Filling:\enspace\enspace 11251 (39.82\%)\\
Gaps Successfully Integrated:\enspace\enspace\enspace 8329\\[0.4em]
Features Passing QC Filter:\enspace\enspace\enspace\enspace\enspace\enspace\enspace 827 / 7064 (11.7\%)
\end{routput}

Correspondence grouping generates $7{,}064$ consensus metabolite features spanning $m/z = 200.1000$ to $600.0000$ Da and retention time $2{,}518.1$ to $4{,}482.6$ seconds. Gap filling recovers $8{,}329$ missing data points, reducing matrix sparsity from $69.30\%$ ($19{,}580$ missing cells) down to $39.82\%$ ($11{,}251$ missing cells). Applying the reproducibility mask ($\text{Det} \ge 50\%$, $\text{CV} < 30\%$) filters out analytical noise, retaining 827 high-integrity features ($11.7\%$) for downstream differential abundance modeling.

\FloatBarrier
\section{Worked Exercises}
\label{sec:worked_exercises}

\subsection*{Exercise 17.1: CentWave Mass Accuracy Parameterization}
\noindent\textbf{Problem Statement:}\enspace
Explain how mass accuracy specifications differ between high-resolution Orbitraps and time-of-flight (TOF) instruments, and state the appropriate \texttt{ppm} parameter setting in \texttt{CentWaveParam} for each platform.

\noindent\textbf{Solution:}\enspace
Orbitrap mass spectrometers achieve resolving powers $R \ge 60{,}000$ and sub-ppm mass accuracy, requiring narrow tolerance windows of $\text{ppm} = 5.0$ to $10.0$. In contrast, standard quadrupole time-of-flight (Q-TOF) instruments operate at $R \approx 10{,}000$ to $30{,}000$ with typical drift of $10$ to $25$ ppm, requiring $\text{ppm} = 25.0$ to $35.0$. 

Setting $\text{ppm} = 30$ on Orbitrap data creates wide regions of interest that merge near-isobaric metabolite ions. Conversely, setting $\text{ppm} = 5$ on TOF data causes the ROI builder to fragment single chromatographic peaks into discontinuous slices, drastically reducing detected peak counts.

\subsection*{Exercise 17.2: Alignment vs Correspondence Mechanics}
\noindent\textbf{Problem Statement:}\enspace
Distinguish the analytical purpose of retention time alignment (\texttt{adjustRtime}) from feature correspondence (\texttt{groupChromPeaks}). Can correspondence grouping run successfully without prior alignment?

\noindent\textbf{Solution:}\enspace
\begin{enumerate}
\item \textbf{Alignment (\texttt{adjustRtime})}: Modifies scan-level retention time coordinates across entire chromatographic runs to correct non-linear physical drift caused by column degradation or mobile phase changes.
\item \textbf{Correspondence (\texttt{groupChromPeaks})}: Matches individual chromatographic peaks across aligned samples into a 2D consensus feature matrix based on Gaussian kernel density estimation.
\end{enumerate}

\noindent\textbf{Execution Feasibility:}\enspace
Correspondence can technically execute without alignment if the density bandwidth $\text{bw}$ is set very wide. However, when run-to-run drift exceeds the physical peak width, a wide bandwidth merges distinct neighboring chromatographic isomers. Conversely, a narrow bandwidth without prior alignment causes the identical chemical metabolite in two runs to split into two separate consensus rows. Alignment is mandatory before correspondence grouping.

\subsection*{Exercise 17.3: Gap Filling Analytical Value}
\noindent\textbf{Problem Statement:}\enspace
In the \texttt{faahKO} cohort analysis, gap filling with \texttt{fillChromPeaks()} reduced matrix missingness from $69.30\%$ to $39.82\%$, recovering $8{,}329$ peak areas. Explain the analytical mechanism of gap filling and describe why it prevents statistical bias in differential testing.

\noindent\textbf{Solution:}\enspace
Peak detection algorithms require peaks to exceed a signal-to-noise ratio threshold ($\text{snthresh} \ge 10$) and scan count threshold ($\text{prefilter}$). When a metabolite is present at low concentration in one biological condition, its elution profile may have $\text{S/N} = 6$, causing the peak picker to register an NA value.

Gap filling identifies the median $m/z$ and retention time bounding box $[m/z_{\min}, m/z_{\max}] \times [\text{RT}_{\min}, \text{RT}_{\max}]$ defined by samples where the peak was detected. It returns to the raw spectral time-series and integrates baseline ion current directly within that bounding box. This prevents artificial truncation and preserves low-level biological signals in linear modeling.

\subsection*{Exercise 17.4: Quality Control Reproducibility Filtering}
\noindent\textbf{Problem Statement:}\enspace
State the theoretical justification for removing features with coefficient of variation $\text{CV} > 30\%$ prior to differential abundance testing.

\noindent\textbf{Solution:}\enspace
In clinical metabolomics, technical variation is quantified across pooled quality control (QC) replicates. A feature exhibiting technical $\text{CV} > 30\%$ has analytical measurement uncertainty that exceeds the typical biological effect size ($\Delta \mu \approx 20\%$--$30\%$). 

Retaining high-CV features in differential linear models inflates multiple hypothesis testing penalties ($P_{\text{adj}}$) through Benjamini-Hochberg adjustment without providing reproducible biological discoveries. Filtering out features with technical $\text{CV} \ge 30\%$ removes analytical noise, enriching the candidate pool for genuine biological signals.

\FloatBarrier
\begin{thebibliography}{10}

\bibitem{smith2006xcms}
Smith, C.A., Want, E.J., O'Maille, G., Abagyan, R., Siuzdak, G.
\newblock {XCMS}: processing mass spectrometry data for metabolite profiling using nonlinear peak alignment, matching, and identification.
\newblock \emph{Analytical Chemistry}, 78(3):779--787, 2006.
\newblock \url{https://doi.org/10.1021/ac051437y}.

\bibitem{tautenhahn2008centwave}
Tautenhahn, R., B\"ottcher, C., Neumann, S.
\newblock Highly sensitive feature detection for high resolution {LC/MS}.
\newblock \emph{BMC Bioinformatics}, 9:504, 2008.
\newblock \url{https://doi.org/10.1186/1471-2105-9-504}.

\bibitem{prince2006obiwarp}
Prince, J.T., Marcotte, E.M.
\newblock Chromatographic alignment of {ESI-LC-MS} proteomics data by ordered bijective interpolated warping.
\newblock \emph{Analytical Chemistry}, 78(17):6140--6152, 2006.
\newblock \url{https://doi.org/10.1021/ac060595t}.

\bibitem{rainer2022msexperiment}
Rainer, J., Vicini, A., Salzer, L., Stanstrup, J., Badia, J.M., Neumann, S., Stravs, M.A., Verri Hernandes, V., Gatto, L.
\newblock A modular and expandable ecosystem for metabolomics data annotation in {R}.
\newblock \emph{Metabolites}, 12(2):173, 2022.
\newblock \url{https://doi.org/10.3390/metabo12020173}.

\bibitem{kuhl2012camera}
Kuhl, C., Tautenhahn, R., B\"ottcher, C., Larson, T.R., Neumann, S.
\newblock {CAMERA}: an integrated strategy for compound spectra extraction and annotation of liquid chromatography/mass spectrometry metabolomics data sets.
\newblock \emph{Analytical Chemistry}, 84(1):283--289, 2012.
\newblock \url{https://doi.org/10.1021/ac202450g}.

\bibitem{broadhurst2018guidelines}
Broadhurst, D., Goodacre, R., Reinke, S.N., et~al.
\newblock Guidelines and considerations for the use of system suitability and quality control samples in mass spectrometry assays applied to untargeted metabolomic workflows.
\newblock \emph{Metabolomics}, 14(6):72, 2018.
\newblock \url{https://doi.org/10.1007/s11306-018-1367-3}.

\end{thebibliography}

\end{document}
